\section{Implementation and Evaluation Methodology}
\label{sec:setup}

\subsection{Prototype and Testbed}
\label{subsec:implementation}

The reference implementation is a Cosmos SDK application, \texttt{engramd},
on a forked CometBFT core. The application mounts no \texttt{x/auth},
\texttt{x/bank}, or \texttt{x/staking}: transactions are unsigned envelopes
and the validator set comes from the genesis file. Three modules hold the
protocol logic: \texttt{x/sovereignty} (the FSM, the sensors, the circuit
breaker, and the ABCI++ hooks), \texttt{x/da} (the Celestia receipt and its
verification), and \texttt{x/anchor} (the Bitcoin settlement receipt and its
re-verification). The ABCI++ hooks follow the pipeline of
Section~\ref{subsec:architecture}: only \texttt{ExtendVote} reads live
sensors, and \texttt{PreBlocker} commits the FSM transition from the agreed
certificate.

The fork differs from upstream CometBFT in four ways: an $f{+}1$ timeout
round-skip, per-peer round-trip-time measurement and a peer health monitor
that feed the P2P sensor, per-peer rate limiting and bounded round acceptance
for timeout messages, and the proposal round passed to the application
through \texttt{PrepareProposal} and \texttt{ProcessProposal}.

Bitcoin and Celestia are live dependencies, deployed as described in
Section~\ref{subsec:testbed}. The recovery proof is a Noir circuit proved with
Barretenberg (Section~\ref{subsec:reanchoring}); a prover process watches for
a \textsc{Sovereign} backlog, builds the proof from committed headers, and
submits it as a transaction that every validator verifies during block
execution. Section~\ref{subsec:environment} lists the versions.

\textbf{Testbed.}\label{subsec:testbed}

The testbed runs the reference implementation of Section~\ref{subsec:implementation}
as seven validators with equal voting power ($n{=}7$, $f{=}2$), each in its
own container. Validator gossip uses a
dedicated full mesh of pairwise links, separate from the network that
carries Bitcoin, Celestia, and RPC traffic; this separation lets us degrade
peripheral connectivity without partitioning the core consensus network.
The peripheral layers are real software, not simulators: a two-node Bitcoin
regtest network mined continuously at roughly one block per 20\,s (thirty times faster than Bitcoin's 600\,s average; the thresholds are in Bitcoin blocks, so detection times measured here convert to mainnet as in Section~\ref{sec:discussion}), and a
Celestia network with seven bridge nodes, one per validator, so that each
validator forms its own DA view. The single-actor attacks of Appendix~\ref{app:attacks} and the $f{=}2$ collusion
suite of Section~\ref{subsec:rq1} ran on the seven-validator cluster.

The Celestia setup is described in Appendix~\ref{app:extended}.

\textbf{Environment and versions.}\label{subsec:environment}

Most results come from Testbed~A: the seven-validator Docker
deployment above on a single host. Unless stated otherwise, the live
tables and figures of Section~\ref{sec:evaluation} are Testbed~A measurements; the
in-process simulations, the microbenchmarks and the offline prover measurements of
Appendix~\ref{app:aggregation} are labelled as such.
\iftestbedB A second environment, Testbed~B, serves the experiments that depend on real
wide-area conditions: seven validators on seven virtual machines in three regions
(us-east-2, eu-west-1 and ap-southeast-1), the Bitcoin regtest node and
\texttt{celestia-app} alone on a machine in a fourth region (us-west-2), and the
prover host with its worker machines in a fifth (us-east-1). Round trips between
regions are 46--207\,ms, and the peer-latency bound of the health sensors is raised to 500\,ms
(Table~\ref{tab:testbedB}). Results from the two environments are never pooled into
one number.\fi{}
Table~\ref{tab:environment} in Appendix~\ref{app:setup} lists the configuration and the versions.


\subsection{Baselines}
\label{subsec:baselines}

\textbf{Strict receipt-gated baseline (primary, RQ2).} The same binary with the
FSM pinned to \textsc{Anchored} (the command-line mode is called
\texttt{hard-gated}). Condition (3) of the validity predicate
(Section~\ref{subsec:proposal}) is enforced at every height, so a proposal is
valid only if it satisfies the agreement checks, carries a Bitcoin receipt within
tolerance, and carries a DA receipt within tolerance:
$V_h(B) = V_{\mathrm{BFT}}(B) \wedge R_{\mathrm{BTC}}(B,h) \wedge R_{\mathrm{DA}}(B,h)$,
with no other validity path while a receipt predicate cannot be satisfied. It is
the behavior of a chain that makes settlement and DA inclusion part of block
validity, it is a construct of this paper and not a published system, and it
isolates the effect of Engram's adaptive relaxation. It is a runtime mode of
the same binary (\texttt{ENGRAM\_MODE=hard-gated}); Section~\ref{subsec:rq2} runs
the RQ2 fault schedule against it. A chain in this mode cannot start from genesis, because the first block needs a DA receipt that does not exist yet; every trial therefore starts the cluster as Engram, waits until all validators report \textsc{Anchored} with healthy sensors, waits until the Bitcoin height gap is at most three, and restarts the validators in the receipt-gated mode with their state kept, with the Bitcoin miner paused during the restart.

The definition of the out-of-band baseline is in Appendix~\ref{app:extended}.

\textbf{Soft baseline (RQ2).} The real FSM runs, but withdrawals are never
locked in \textsc{Sovereign} or \textsc{Recovering}
(\texttt{ENGRAM\_MODE=soft}). It keeps liveness and drops the circuit
breaker, which isolates what the withdrawal lock contributes.

The vanilla control and the ablation baselines are defined in Appendix~\ref{app:extended}.

\subsection{Fault Model, Schedules and Parameters}
The fault schedules and the fault model are in Appendix~\ref{app:extended}.

\textbf{FSM parameters.}\label{subsec:params}

The parameters are fixed in the genesis file, so every validator holds identical
values; Table~\ref{tab:fsm-params} in Appendix~\ref{app:setup} lists them. The live
experiments use a Bitcoin gap of 5 blocks for \textsc{Suspicious} and 8 for
\textsc{Sovereign}, a DA gap of 30 blocks, a finality depth
$K_{\text{DEEP\_FINALITY}}{=}2$, hysteresis waits of 2 blocks, a down threshold
$k_{\downarrow}{=}2$, a \textsc{Suspicious} escalation bound of 24 blocks, and a
recovery gate that accepts only a verified proof with an unproven tail of at most
128 blocks, unless a section states an override. Sensors are sampled once per
block, when a validator signs its vote extension; the measurement scripts poll
each validator's RPC every 3\,s (2\,s during the packet-loss cycles of C4). Blocks
arrive about once per second.


\subsection{Workloads, Metrics and Statistics}
\label{subsec:metrics}

The synthetic workloads are described in Appendix~\ref{app:extended}.

\textbf{Terms.} $n$ is the number of validators ($n{=}7$ on the main testbed). A \emph{run} is one execution of a
scenario; unless stated otherwise it starts from \textsc{Anchored} on a freshly
deployed chain, so runs are independent deployments, and intervals over runs
are intervals over deployments. A \emph{repeat} is a further measurement on the
chain of an earlier one (the repeats of Table~\ref{tab:overhead} share a
deployment, so their interval describes one chain). An \emph{arm} is one
system configuration that a table compares with the others. The number of runs
is written out wherever it matters. Every 95\% interval over runs uses the Student $t$
critical value for the number of runs (for example 2.78 for five), not 1.96.

Let a run record the FSM state $s_b$ of every block $b$ in a window $W$.
\begin{itemize}
  \item \emph{Anchored uptime} $=|\{b\in W : s_b=\textsc{Anchored}\}|/|W|$.
  In-process, $W$ is the whole recorded timeline: for the noisy
  environments, the 100 noisy blocks plus the setup blocks that precede them
  (two for the \textsc{Recovering}$\to$\textsc{Anchored} sweep, one for the
  others). The no-noise control of that sweep ends at the first
  \textsc{Anchored} block, so its window is only the setup blocks plus the
  wait, and its uptime falls as the wait grows. Live, $W$ is the sequence of
  RPC samples of one validator over a 300\,s window.
  \item \emph{Flapping count} is the number of transitions that immediately
  reverse the preceding transition, for example
  \textsc{Anchored}$\to$\textsc{Suspicious} followed by
  \textsc{Suspicious}$\to$\textsc{Anchored}, over the same $W$. Live runs
  detect reversals between consecutive samples, so a reversal faster than
  the polling interval is not seen.
  \item \emph{Absorption rate} is $a/(a+d)$, where $a$ is the number of
  disturbed blocks that begin in \textsc{Anchored} and leave the state
  unchanged, and $d$ is the number that move it to \textsc{Suspicious}.
  It is defined only on the \textsc{Anchored}$\to$\textsc{Suspicious} and
  \textsc{Suspicious}$\to$\textsc{Anchored} sweeps.
  \item \emph{Blocks/s} $=(h_{\max}-h_{\min})/T$ over a measurement window of
  length $T$ (60\,s after an 8\,s warm-up in Section~\ref{subsec:rq4}).
  \emph{Steady blocks/s} $=1/\bar\imath$, where $\bar\imath$ is the mean
  block interval. \emph{Gap fraction} $=1-n\bar\imath/T$, where $n$ is the
  number of observed intervals: the share of the window not accounted for
  by intervals at the steady pace, that is, time spent in heights that took
  longer than one interval.
\end{itemize}
